\documentclass[a4paper,12pt]{article}

\usepackage[spanish,es-noquoting]{babel}
\usepackage[utf8]{inputenc}
\usepackage[T1]{fontenc}
\usepackage{amsmath, amssymb}
\usepackage{enumitem}
\usepackage{tikz}
\usetikzlibrary{arrows.meta,calc,patterns}
\usepackage[margin=2cm]{geometry}
\usepackage[hidelinks]{hyperref}
\usepackage{parskip}

\newcommand{\dd}{\,\mathrm{d}}

\begin{document}

\begin{center}
    \Large\textbf{Problema parcial 4}\\
    \normalsize Física Matemática I
\end{center}

\tableofcontents

\section{Enunciado}

Considere una esfera de radio $R$ con densidad
\[
\rho(r,\theta,\varphi)=
\begin{cases}
\rho_0\left(\dfrac{r}{R}\right)e^{\alpha r/R}\sin^2\theta\cos(2\varphi), & 0\le r\le R,\\
0, & r>R,
\end{cases}
\]
donde $\alpha$ y $\rho_0$ son constantes. El potencial satisface la ecuación de Poisson
$\nabla^2\phi(\vec r)=-\rho(\vec r)/\epsilon_0$, cuya solución mediante la función de Green es
\begin{equation}
\label{eq:green-potential}
\phi(\vec r)=\frac{1}{4\pi\epsilon_0}\int \frac{\rho(\vec r\,')}{\lVert \vec r-\vec r\,'\rVert}\dd^3r'.
\end{equation}

Para la función de Green se usa la expansión
\[
\frac{1}{\lVert \vec r-\vec r\,'\rVert}
=\begin{cases}
\displaystyle \sum_{\ell=0}^{\infty}\frac{r^\ell}{r'^{\ell+1}}P_\ell(\cos\gamma), & r<r',\\[0.5em]
\displaystyle \sum_{\ell=0}^{\infty}\frac{r'^\ell}{r^{\ell+1}}P_\ell(\cos\gamma), & r>r',
\end{cases}
\]
donde $\gamma$ es el ángulo entre $\vec r$ y $\vec r\,'$. Además, por el teorema de adición,
\begin{equation}
\label{eq:addition-theorem}
P_\ell(\cos\gamma)=\frac{4\pi}{2\ell+1}\sum_{m=-\ell}^{\ell}
Y_{\ell m}^*(\theta,\varphi)Y_{\ell m}(\theta',\varphi').
\end{equation}

Se pide:
\begin{enumerate}[label=\alph*)]
    \item Encontrar el potencial eléctrico en el interior de la esfera, $0\le r\le R$.
    \item Determinar el campo eléctrico para $0\le r\le R$.
    \item Encontrar el campo eléctrico en el exterior, $r>R$.
    \item Dibujar la magnitud del campo eléctrico y del potencial.
\end{enumerate}

\section{Potencial eléctrico en el interior}

\begin{center}
\begin{tikzpicture}[scale=1.35, line cap=round, line join=round, >=Latex]
    \shade[inner color=orange!18, outer color=orange!4, draw=none] (0,0) circle (0.88);
    \shade[inner color=blue!4, outer color=blue!13, draw=none, even odd rule]
        (0,0) circle (1.62) (0,0) circle (0.88);
    \draw[blue!65!black, very thick] (0,0) circle (1.62);
    \draw[orange!85!black, very thick, dashed] (0,0) circle (0.88);
    \draw[gray!40, dashed] (-1.62,0) -- (1.62,0);
    \draw[gray!40, dashed] (0,-1.62) -- (0,1.62);
    \draw[->, orange!85!black, very thick] (0,0) -- (0.88,0)
        node[midway, above] {$r$};
    \draw[->, blue!65!black, very thick] (0,0) -- ({1.62*cos(38)},{1.62*sin(38)})
        node[midway, above left] {$R$};
    \draw[orange!85!black, thick] (0.55,-0.55) -- (2.05,-1.05);
    \draw[blue!65!black, thick] (-1.18,-0.95) -- (-2.3,-1.28);
    \fill (0,0) circle (1.1pt) node[below left] {$O$};
    \node[font=\bfseries] at (0,1.95) {Separación radial de la integral};
    \node[orange!85!black, align=center, fill=white, rounded corners=2pt, inner sep=3pt]
        at (2.75,-1.12) {Región I\\$0\le r'\le r$};
    \node[blue!65!black, align=center, fill=white, rounded corners=2pt, inner sep=3pt]
        at (-3.0,-1.35) {Región II\\$r\le r'\le R$};
\end{tikzpicture}
\end{center}

La dependencia angular de la densidad es proporcional a
$\sin^2\theta\cos(2\varphi)$. Como
\[
Y_{\ell m}(\theta,\varphi)=\sqrt{\frac{2\ell+1}{4\pi}\frac{(\ell-m)!}{(\ell+m)!}}
P_\ell^m(\cos\theta)e^{im\varphi},
\]
los términos relevantes son $\ell=2$ y $m=\pm2$. En la expansión multipolar, el índice $\ell$ clasifica el tipo de simetría angular: $\ell=0$ corresponde al término monopolar, $\ell=1$ al dipolar, $\ell=2$ al cuadrupolar, $\ell=3$ al octupolar, y así sucesivamente. En este caso, $\ell=2$ significa que la distribución tiene carácter cuadrupolar: presenta cuatro lóbulos angulares alternados, con regiones de signo opuesto, en lugar de una simetría puramente radial o dipolar. El índice $m=\pm2$ fija la variación azimutal proporcional a $\cos(2\varphi)$.

En particular,
\[
Y_{2,2}(\theta,\varphi)=\sqrt{\frac{5}{96\pi}}\,3\sin^2\theta\,e^{2i\varphi},
\qquad
Y_{2,-2}(\theta,\varphi)=\sqrt{\frac{5}{96\pi}}\,3\sin^2\theta\,e^{-2i\varphi}.
\]
Por tanto,
\[
\sin^2\theta\cos(2\varphi)
=\sqrt{\frac{8\pi}{15}}\left[Y_{2,2}(\theta,\varphi)+Y_{2,-2}(\theta,\varphi)\right].
\]
Así, la densidad queda escrita como
\begin{equation}
\label{eq:density-harmonics}
\rho(r',\theta',\varphi')=\rho_0\left(\frac{r'}{R}\right)e^{\alpha r'/R}
\sqrt{\frac{8\pi}{15}}\left[Y_{2,2}(\theta',\varphi')+Y_{2,-2}(\theta',\varphi')\right].
\end{equation}

Para un punto interior, $0\le r\le R$, la integral radial debe separarse en dos regiones:
$0\le r'\le r$, donde $r>r'$, y $r\le r'\le R$, donde $r<r'$. Usando \eqref{eq:green-potential},
\[
\phi_{\mathrm{int}}(r,\theta,\varphi)
=\frac{1}{4\pi\epsilon_0}\int_0^R\int_{\Omega'}
\frac{\rho(r',\theta',\varphi')}{\lVert\vec r-\vec r\,'\rVert}r'^2\dd r'\dd\Omega'.
\]
El elemento de volumen en coordenadas esféricas es
\[
\dd^3r'=r'^2\dd r'\dd\Omega',
\qquad
\int_{\Omega'}\dd\Omega'=\int_0^\pi\int_0^{2\pi}\sin\theta'\dd\varphi'\dd\theta'.
\]

En la Región I se reemplaza
\[
\frac{1}{\lVert\vec r-\vec r\,'\rVert}
=\sum_{\ell=0}^{\infty}\frac{r'^\ell}{r^{\ell+1}}P_\ell(\cos\gamma),
\]
y en la Región II,
\[
\frac{1}{\lVert\vec r-\vec r\,'\rVert}
=\sum_{\ell=0}^{\infty}\frac{r^\ell}{r'^{\ell+1}}P_\ell(\cos\gamma).
\]
Al sustituir \eqref{eq:density-harmonics} y \eqref{eq:addition-theorem}, resulta
\begin{align*}
\phi_{\mathrm{int}}(r,\theta,\varphi)
&=\frac{\rho_0}{4\pi\epsilon_0}\sqrt{\frac{8\pi}{15}}
\sum_{\ell=0}^{\infty}\sum_{m=-\ell}^{\ell}\frac{4\pi}{2\ell+1}Y_{\ell m}^*(\theta,\varphi) \\
&\quad \times \left[
\int_0^r\frac{r'^\ell}{r^{\ell+1}}\frac{r'}{R}r'^2e^{\alpha r'/R}\dd r'
\int_{\Omega'}Y_{\ell m}(\theta',\varphi')\left(Y_{2,2}+Y_{2,-2}\right)\dd\Omega'
\right.\\
&\quad\left.+
\int_r^R\frac{r^\ell}{r'^{\ell+1}}\frac{r'}{R}r'^2e^{\alpha r'/R}\dd r'
\int_{\Omega'}Y_{\ell m}(\theta',\varphi')\left(Y_{2,2}+Y_{2,-2}\right)\dd\Omega'
\right].
\end{align*}

Por ortogonalidad angular solo sobreviven $\ell=2$ y $m=\pm2$. Entonces
\[
\phi_{\mathrm{int}}(r,\theta,\varphi)
=\frac{\rho_0}{4\pi\epsilon_0}\sqrt{\frac{8\pi}{15}}\frac{4\pi}{5}
\left[\frac{1}{Rr^3}\int_0^r r'^5e^{\alpha r'/R}\dd r'
+\frac{r^2}{R}\int_r^R e^{\alpha r'/R}\dd r'\right]
\left(Y_{2,2}^*+Y_{2,-2}^*\right).
\]
Definimos
\[
I_1(r)=\int_0^r r'^5e^{\alpha r'/R}\dd r',
\qquad
I_2(r)=\int_r^R e^{\alpha r'/R}\dd r'.
\]
Estas integrales son
\begin{align*}
I_1(r)
&=e^{\alpha r/R}\left[\frac{Rr^5}{\alpha}-\frac{5R^2r^4}{\alpha^2}
+\frac{20R^3r^3}{\alpha^3}-\frac{60R^4r^2}{\alpha^4}
+\frac{120R^5r}{\alpha^5}-\frac{720R^6}{\alpha^6}\right]
+\frac{720R^6}{\alpha^6},\\
I_2(r)&=\frac{R}{\alpha}\left[e^\alpha-e^{\alpha r/R}\right].
\end{align*}

Usando de nuevo $Y_{2,2}^*+Y_{2,-2}^*=\sqrt{15/(8\pi)}\sin^2\theta\cos(2\varphi)$, se obtiene
\begin{equation}
\label{eq:phi-int}
\boxed{
\begin{aligned}
\phi_{\mathrm{int}}(r,\theta,\varphi)
&=\frac{\rho_0}{5\epsilon_0R}\sin^2\theta\cos(2\varphi)
\left[\frac{I_1(r)}{r^3}+r^2I_2(r)\right].
\end{aligned}}
\end{equation}

\section{Campo eléctrico en el interior}

\begin{center}
\begin{tikzpicture}[scale=1.35, line cap=round, line join=round, >=Latex]
    \shade[ball color=blue!8, opacity=0.55] (0,0) circle (1.55);
    \draw[blue!60!black, very thick] (0,0) circle (1.55);
    \draw[->, gray!70] (-1.95,0) -- (2.05,0) node[right] {$x$};
    \draw[->, gray!70] (0,-1.95) -- (0,2.05) node[above] {$y$};
    \foreach \x/\y/\u/\v in {-1.05/0.95/-0.52/-0.16, 1.05/0.95/0.52/-0.16,
        -1.05/-0.95/-0.45/0.22, 1.05/-0.95/0.45/0.22,
        -0.42/0.58/-0.34/-0.30, 0.42/0.58/0.34/-0.30,
        -0.42/-0.58/-0.30/0.34, 0.42/-0.58/0.30/0.34,
        0/1.28/0/-0.48, 0/-1.28/0/0.48}{
        \draw[->, red!70!black, very thick] (\x,\y) -- ++(\u,\v);
    }
    \node[blue!60!black] at (1.78,0.25) {$r=R$};
    \node[font=\bfseries, align=center] at (0,-2.35) {Campo interior\\$\vec E_{\mathrm{int}}=-\nabla\phi_{\mathrm{int}}$};
\end{tikzpicture}
\end{center}

El campo eléctrico se obtiene a partir de
\[
\vec E=-\vec\nabla\phi
=-\left(\frac{\partial\phi}{\partial r}\hat r
+\frac{1}{r}\frac{\partial\phi}{\partial\theta}\hat\theta
+\frac{1}{r\sin\theta}\frac{\partial\phi}{\partial\varphi}\hat\varphi\right).
\]
De \eqref{eq:phi-int}, escribimos
\[
\phi_{\mathrm{int}}(r,\theta,\varphi)=f(r)\sin^2\theta\cos(2\varphi),
\qquad
f(r)=\frac{\rho_0}{5\epsilon_0R}\left[\frac{I_1(r)}{r^3}+r^2I_2(r)\right].
\]
Al derivar,
\[
f'(r)=\frac{\rho_0}{5\epsilon_0R}\left[-\frac{3I_1(r)}{r^4}+2rI_2(r)\right],
\]
porque los términos $r^2e^{\alpha r/R}$ que provienen de $I_1'(r)$ y de $I_2'(r)$ se cancelan.

Las componentes interiores son
\begin{align*}
E_r^{\mathrm{int}}&=-f'(r)\sin^2\theta\cos(2\varphi),\\
E_\theta^{\mathrm{int}}&=-\frac{f(r)}{r}\sin(2\theta)\cos(2\varphi),\\
E_\varphi^{\mathrm{int}}&=\frac{2f(r)}{r}\sin\theta\sin(2\varphi).
\end{align*}
Por tanto,
\begin{equation}
\label{eq:e-int}
\boxed{
\vec E_{\mathrm{int}}
=-f'(r)\sin^2\theta\cos(2\varphi)\hat r
-\frac{f(r)}{r}\sin(2\theta)\cos(2\varphi)\hat\theta
+\frac{2f(r)}{r}\sin\theta\sin(2\varphi)\hat\varphi.}
\end{equation}

\section{Campo eléctrico en el exterior}

\begin{center}
\begin{tikzpicture}[scale=1.35, line cap=round, line join=round, >=Latex]
    \shade[ball color=blue!10, opacity=0.6] (0,0) circle (0.95);
    \draw[blue!60!black, very thick] (0,0) circle (0.95);
    \draw[gray!45, dashed, thick] (0,0) circle (1.65);
    \draw[gray!30, dashed] (0,0) circle (2.2);
    \foreach \ang/\len in {20/0.62,55/0.45,100/0.58,145/0.45,200/0.62,235/0.45,280/0.58,325/0.45}{
        \draw[->, red!70!black, very thick]
            ({1.30*cos(\ang)},{1.30*sin(\ang)}) --
            ({(1.30+\len)*cos(\ang)},{(1.30+\len)*sin(\ang)});
    }
    \draw[->, blue!60!black, very thick] (0,0) -- (0.95,0) node[midway, above] {$R$};
    \node[gray!70!black] at (1.88,0.27) {$r>R$};
    \node[font=\bfseries, align=center] at (0,-2.55) {Campo exterior\\decaimiento cuadrupolar $\propto r^{-4}$};
\end{tikzpicture}
\end{center}

Para $r>R$, todo punto de la fuente cumple $r'>0$ y $r>r'$, de modo que
\[
\frac{1}{\lVert\vec r-\vec r\,'\rVert}
=\sum_{\ell=0}^{\infty}\frac{r'^\ell}{r^{\ell+1}}P_\ell(\cos\gamma).
\]
Usando \eqref{eq:density-harmonics} y \eqref{eq:addition-theorem} en \eqref{eq:green-potential},
\begin{align*}
\phi_{\mathrm{ext}}(r,\theta,\varphi)
&=\frac{\rho_0}{4\pi\epsilon_0}\sqrt{\frac{8\pi}{15}}
\sum_{\ell=0}^{\infty}\sum_{m=-\ell}^{\ell}\frac{4\pi}{2\ell+1}Y_{\ell m}^*(\theta,\varphi)\\
&\quad\times \int_0^R\frac{r'^\ell}{r^{\ell+1}}\frac{r'}{R}r'^2e^{\alpha r'/R}\dd r'
\int_{\Omega'}Y_{\ell m}(\theta',\varphi')\left(Y_{2,2}+Y_{2,-2}\right)\dd\Omega'.
\end{align*}
Nuevamente, la integral angular selecciona solo $\ell=2$ y $m=\pm2$. Así,
\[
\phi_{\mathrm{ext}}(r,\theta,\varphi)
=\frac{\rho_0}{4\pi\epsilon_0}\sqrt{\frac{8\pi}{15}}\frac{4\pi}{5}
\frac{I_1(R)}{Rr^3}\sqrt{\frac{15}{8\pi}}\sin^2\theta\cos(2\varphi).
\]
Como
\[
I_1(R)=e^\alpha R^6\left[\frac{1}{\alpha}-\frac{5}{\alpha^2}+\frac{20}{\alpha^3}
-\frac{60}{\alpha^4}+\frac{120}{\alpha^5}-\frac{720}{\alpha^6}\right]
+\frac{720R^6}{\alpha^6},
\]
el potencial exterior queda
\begin{equation}
\label{eq:phi-ext}
\boxed{
\phi_{\mathrm{ext}}(r,\theta,\varphi)
=\frac{\rho_0R^5}{5\epsilon_0\alpha^6r^3}
\left[e^\alpha(\alpha^5-5\alpha^4+20\alpha^3-60\alpha^2+120\alpha-720)+720\right]
\sin^2\theta\cos(2\varphi).}
\end{equation}

Para calcular el campo, definimos el factor común
\begin{equation}
\label{eq:q-ext}
Q=\frac{\rho_0R^5}{5\epsilon_0\alpha^6}
\left[e^\alpha(\alpha^5-5\alpha^4+20\alpha^3-60\alpha^2+120\alpha-720)+720\right].
\end{equation}
Entonces \eqref{eq:phi-ext} se escribe como
\[
\phi_{\mathrm{ext}}(r,\theta,\varphi)=\frac{Q}{r^3}\sin^2\theta\cos(2\varphi).
\]
Derivando en coordenadas esféricas,
\begin{align*}
E_r^{\mathrm{ext}}
&=-\frac{\partial}{\partial r}\left[\frac{Q}{r^3}\sin^2\theta\cos(2\varphi)\right]
=\frac{3Q}{r^4}\sin^2\theta\cos(2\varphi),\\
E_\theta^{\mathrm{ext}}
&=-\frac{1}{r}\frac{\partial}{\partial\theta}\left[\frac{Q}{r^3}\sin^2\theta\cos(2\varphi)\right]
=-\frac{Q}{r^4}\sin(2\theta)\cos(2\varphi),\\
E_\varphi^{\mathrm{ext}}
&=-\frac{1}{r\sin\theta}\frac{\partial}{\partial\varphi}\left[\frac{Q}{r^3}\sin^2\theta\cos(2\varphi)\right]
=\frac{2Q}{r^4}\sin\theta\sin(2\varphi).
\end{align*}
Por tanto,
\begin{equation}
\label{eq:e-ext}
\boxed{
\vec E_{\mathrm{ext}}
=\frac{Q}{r^4}\left[3\sin^2\theta\cos(2\varphi)\hat r
-\sin(2\theta)\cos(2\varphi)\hat\theta
+2\sin\theta\sin(2\varphi)\hat\varphi\right].}
\end{equation}

\section{Indicaciones para graficar}

Para construir una gráfica representativa en el plano $xy$, conviene fijar valores adimensionales sencillos. Una elección adecuada es
\[
R=1,\qquad \alpha=1,\qquad \rho_0=1,\qquad \epsilon_0=1.
\]
Estos valores no buscan representar un material específico, sino normalizar el problema para visualizar claramente la estructura angular del potencial y del campo. Con $R=1$, el interior corresponde al disco $x^2+y^2\le1$ y el exterior a $x^2+y^2>1$.

Como el plot será únicamente en el plano $xy$, se debe tomar $z=0$. En coordenadas esféricas esto corresponde a $\theta=\pi/2$, mientras que
\[
r=\sqrt{x^2+y^2},\qquad \varphi=\operatorname{atan2}(y,x).
\]
En cada punto de una malla rectangular del plano $xy$, se evalúa si $r\le R$ o si $r>R$. Para los puntos interiores se usa \eqref{eq:phi-int} y \eqref{eq:e-int}; para los puntos exteriores se usa \eqref{eq:phi-ext} y \eqref{eq:e-ext}.

Para graficar el potencial se puede representar directamente $\phi(x,y)$. Para graficar la magnitud del campo eléctrico se debe calcular
\[
\lVert\vec E\rVert=\sqrt{E_r^2+E_\theta^2+E_\varphi^2},
\]
usando las componentes interiores o exteriores según corresponda. Si se desea una visualización suave cerca del origen, se recomienda evitar evaluar exactamente en $r=0$ o reemplazar ese punto por un valor límite/nulo apropiado para la escala de la gráfica.

Una buena presentación consiste en hacer dos mapas de color en el plano $xy$: uno para $\phi(x,y)$ y otro para $\lVert\vec E(x,y)\rVert$. También es útil marcar la circunferencia $r=R$ para distinguir visualmente la frontera entre interior y exterior.

\end{document}